problem:  
Let \( G \) be a simply connected **nilpotent Lie group** of dimension \( n \) with Lie algebra \( \mathfrak{g} \). Each left-invariant Riemannian metric \( g \) on \( G \) determines a Ricci operator \( \mathrm{Ric}_g \!:\! \mathfrak{g} \to \mathfrak{g} \). Let \( g(t) \) evolve by **normalized Ricci flow** so that  
\[
\frac{\partial g}{\partial t} = -2 \big( \mathrm{Ric}_g - \tfrac{r_g}{n} g \big),
\]
and let \( \mathrm{Ric}_t \) denote the corresponding time-\(t\) Ricci operator, viewed as a symmetric operator with respect to \( g(t) \).

For a nilpotent \( \mathfrak{g} \), the corresponding normalized **bracket flow** in Lauret’s sense is known to converge, up to automorphisms, to a **limiting bracket** \( \mu_\infty \) which is a nilsoliton bracket when one exists.  In such a case the nilsoliton Ricci operator satisfies  
\[
\mathrm{Ric}_{\infty} = c I + D_\infty,
\]
where \( D_\infty \) is a derivation of \( \mathfrak{g} \) uniquely determined up to scaling and inner derivations.

**Open problem:**  
Let \( g(t) \) be the normalized Ricci flow on a nilpotent \( G \).  Denote by \( \{\lambda_i(t)\}_{i=1}^n \) the ordered eigenvalues of \( \mathrm{Ric}_t \) and by \( E_i(t) \subset \mathfrak{g} \) the associated eigenspaces (each \(E_i(t)\) determined up to orientation).  Consider the Grassmannian topology on the set of subspaces of \( \mathfrak{g} \).

1. (**Eigenflag alignment**)  When the flow converges algebraically to a nilsoliton bracket \( \mu_\infty \), does there exist a canonical filtration  
   \[
   0 \subsetneq V_1 \subsetneq \cdots \subsetneq V_k = \mathfrak{g},
   \]
   determined only by the Lie algebraic data of \( (\mathfrak{g}, D_\infty) \), such that the evolving eigenspaces \( E_i(t) \) become asymptotically tangent to these \( V_j \) in the Grassmannian, i.e.  
   \[
   \mathrm{dist}(E_i(t),V_j)\to 0 \quad \text{for suitable assignments } i\mapsto j?
   \]
   Can such a flag be characterized purely in terms of the **weight-space decomposition** of the derivation \(D_\infty\), independent of the initial metric?

2. (**Spectral detection of non-solitons**)  When \(\mathfrak{g}\) admits no nilsoliton, do the trajectories of the renormalized eigenvalue ratios
   \[
   \rho_i(t) := \frac{\lambda_i(t)-\frac{1}{n}\sum_j \lambda_j(t)}{\big(\sum_j (\lambda_j(t)-\frac{1}{n}\sum_k \lambda_k(t))^2\big)^{1/2}}
   \]
   converge to characteristic limit points determined by the orbit closure of \([\mu]\) under \(GL(n,\mathbb{R})\)?  Equivalently, can one detect the degeneration direction of the bracket flow purely from the limiting spectrum of \(\mathrm{Ric}_t\)?

3. (**Universality of spectral asymptotics**)  Are the asymptotic ratios \(\rho_i(t)/\rho_j(t)\) rational functions of the eigenvalues of the pre‑Einstein derivation of \(\mathfrak{g}\)?  If so, this would endow the Ricci‑flow phase space with a canonical **spectral stratification** unifying geometric and algebraic Ricci data.

Why is it a "good" problem:  
This formulation links the **spectral geometry** of the Ricci operator to **algebraic derivation structures** intrinsic to the nilpotent Lie algebra, focusing on rigorously defined objects—eigenvalue ratios, Grassmannian convergence of eigenspaces, and derivation gradings.  

It asks whether the Ricci flow dynamically reconstructs, from analytical evolution alone, the canonical flag determined by the nilsoliton or pre‑Einstein derivation.  The question is simple to state (tracking eigenstructures of a symmetric operator under an evolution equation) but probes deeply into the intersection of **geometric analysis**, **representation theory of Lie algebras**, and **algebraic dynamics of moment‑map flows**.  

A positive resolution would reveal that the Ricci flow "learns" the internal grading of a Lie algebra through its spectral evolution; a negative one would highlight new types of degenerations invisible to algebraic derivations.  Either outcome would substantially extend our conceptual understanding of homogeneous Ricci dynamics, making this a genuinely “good” and forward‑looking research problem.